\documentclass[11pt]{article}
\usepackage[margin=1in]{geometry}
\usepackage{amsmath,amssymb}
\usepackage{booktabs}
\usepackage{listings}
\usepackage{xcolor}
\usepackage{hyperref}

\lstset{
  language=C++,
  basicstyle=\ttfamily\small,
  keywordstyle=\color{blue},
  commentstyle=\color{gray},
  breaklines=true,
  frame=single
}

\title{Stochastic Supernova Driving in AthenaK:\\Thermal Energy and Radial Momentum Injection}
\author{}
\date{\today}

\begin{document}
\maketitle

\section{Overview}

The supernova (SN) driving module in AthenaK implements stochastic injection of
thermal energy and radial momentum into a hydrodynamic or MHD simulation.  Events
are generated as a Poisson-like process and injected into spherical regions of
radius $r_\mathrm{inj}$.  The two injection channels---thermal energy
$E_\mathrm{inj}$ and radial momentum $P_\mathrm{inj}$---are independently
configurable, allowing pure thermal feedback, pure momentum feedback, or any
combination.

\section{Event Generation}

\subsection{Stochastic Sampling}

Events are drawn from a quasi-Poisson process controlled by a target rate
$\dot{N}_\mathrm{SN}$ (events per unit code time).  At each hydrodynamic
substep with effective timestep $\beta\,\Delta t$, the fractional event
accumulator is updated on rank~0:
\begin{equation}
  A \leftarrow A + \dot{N}_\mathrm{SN}\,\beta\,\Delta t\,,
\end{equation}
where $A$ is carried between steps.  The integer part gives the number of events
for this step:
\begin{equation}
  N_\mathrm{ev} = \lfloor A \rfloor\,, \qquad A \leftarrow A - N_\mathrm{ev}\,.
\end{equation}
This preserves the long-term average rate while distributing events discretely
across timesteps.

\subsection{Spatial Placement}

Each event centre $(x_\mathrm{sn}, y_\mathrm{sn}, z_\mathrm{sn})$ is drawn
uniformly:
\begin{align}
  x_\mathrm{sn} &\sim \mathcal{U}(x_{1,\min},\, x_{1,\max})\,, \\
  y_\mathrm{sn} &\sim \mathcal{U}(x_{2,\min},\, x_{2,\max})\,, \\
  z_\mathrm{sn} &\sim \mathcal{U}(z_\mathrm{band,min},\, z_\mathrm{band,max})\,,
\end{align}
where the $z$-band $[z_\mathrm{band,min},\, z_\mathrm{band,max}]$ is a
user-specified sub-interval of the $x_3$ domain, clipped to the mesh extent.
For lower-dimensional problems, unused coordinates default to the domain
midpoint.

Event coordinates are generated on MPI rank~0 using a 64-bit Mersenne Twister
(\texttt{std::mt19937\_64}) seeded deterministically, then broadcast to all
ranks.

\section{Injection Geometry}

The injection region for each event is a tophat sphere of effective radius
$r_\mathrm{inj}$:
\begin{equation}
  \mathcal{S} = \bigl\{ \mathbf{x} : |\mathbf{x} - \mathbf{x}_\mathrm{sn}|^2 < r_\mathrm{inj}^2 \bigr\}\,.
\end{equation}
A minimum resolved radius is enforced:
\begin{equation}
  r_\mathrm{inj,eff} = \max\!\bigl(r_\mathrm{inj},\; 2\,\Delta x_\mathrm{min}\bigr)\,,
\end{equation}
where $\Delta x_\mathrm{min}$ is the global minimum cell width across all
directions and all MeshBlocks (obtained via \texttt{MPI\_Allreduce}).

Periodic boundary wrapping is applied to each displacement component.  For
example, in the $x_1$ direction with domain length $L_{x_1}$:
\begin{equation}
  \delta x = x - x_\mathrm{sn}\,; \qquad
  \text{if } \delta x > \tfrac{1}{2}L_{x_1},\; \delta x \leftarrow \delta x - L_{x_1}\,; \quad
  \text{if } \delta x < -\tfrac{1}{2}L_{x_1},\; \delta x \leftarrow \delta x + L_{x_1}\,.
\end{equation}

\section{Three-Pass Algorithm}

For each event, three parallel passes over the mesh are performed.

\subsection{Pass 1: Volume Calculation}

A Kokkos parallel reduction accumulates the total volume of cells within the
injection sphere:
\begin{equation}
  V = \sum_{c \,\in\, \mathcal{S}} \Delta x_{1,c}\,\Delta x_{2,c}\,\Delta x_{3,c}\,.
\end{equation}
This is reduced across MPI ranks to obtain the global volume $V$.

\subsection{Pass 2: Density Averaging (Momentum Injection Only)}

When $P_\mathrm{inj} > 0$, a second parallel reduction computes the
volume-weighted mean density within the sphere:
\begin{equation}
  \bar{\rho} = \frac{1}{V} \sum_{c \,\in\, \mathcal{S}} \rho_c \,\Delta x_{1,c}\,\Delta x_{2,c}\,\Delta x_{3,c}\,,
\end{equation}
where $\rho_c = w_0(\texttt{IDN})$ is the primitive density.  The numerator is
reduced across MPI ranks before dividing by $V$.

The uniform radial velocity is then:
\begin{equation}\label{eq:vri}
  v_r = \frac{P_\mathrm{inj}}{\bar{\rho}\,V}\,.
\end{equation}
This ensures that the total injected radial momentum equals
$P_\mathrm{inj}$:
\begin{equation}
  \sum_{c \,\in\, \mathcal{S}} \bar{\rho}\,v_r\,\Delta V_c = \bar{\rho}\,v_r\,V = P_\mathrm{inj}\,.
\end{equation}

This pass is skipped entirely when $P_\mathrm{inj} = 0$.

\subsection{Pass 3: Injection}

For each cell $c \in \mathcal{S}$ with displacement
$\boldsymbol{\delta} = (\delta x,\, \delta y,\, \delta z)$ from the event
centre and $r = |\boldsymbol{\delta}|$:

\paragraph{Thermal energy injection (when $E_\mathrm{inj} > 0$):}
The energy density increment is distributed uniformly over the sphere volume:
\begin{equation}
  \Delta e = \frac{E_\mathrm{inj}}{V}\,,
\end{equation}
and added to the total energy:
\begin{equation}
  u_0(\texttt{IEN})_c \mathrel{+}= \Delta e\,.
\end{equation}

\paragraph{Momentum injection (when $P_\mathrm{inj} > 0$ and $r > 0$):}
Radial momentum is injected along the unit vector
$\hat{\boldsymbol{r}} = \boldsymbol{\delta}/r$:
\begin{align}
  u_0(\texttt{IM1})_c &\mathrel{+}= \bar{\rho}\,v_r\,\frac{\delta x}{r}\,, \\
  u_0(\texttt{IM2})_c &\mathrel{+}= \bar{\rho}\,v_r\,\frac{\delta y}{r}\,, \\
  u_0(\texttt{IM3})_c &\mathrel{+}= \bar{\rho}\,v_r\,\frac{\delta z}{r}\,,
\end{align}
with corresponding kinetic energy:
\begin{equation}
  u_0(\texttt{IEN})_c \mathrel{+}= \tfrac{1}{2}\,\bar{\rho}\,v_r^2\,.
\end{equation}

\paragraph{Centre cell handling ($r = 0$):}
When the event centre coincides exactly with a cell centre, $r = 0$ and the
unit vector $\hat{\boldsymbol{r}}$ is undefined.  Following the approach used
in RAMSES---where randomly placed event centres have zero probability of
coinciding exactly with a cell centre---the momentum injection is simply
skipped for the centre cell.  The cell still receives thermal energy.  A
guard \texttt{r2 > 0} in the GPU kernel prevents any floating-point
exception.

\section{Numerical Requirements: First-Order Flux Correction}

Momentum injection creates steep velocity gradients at the boundary of the
injection sphere.  With a second-order (PLM) spatial reconstruction scheme,
these gradients can cause reconstruction overshoots in adjacent cells, driving
the local internal energy negative and producing floating-point exceptions
(\texttt{NaN}) in the Riemann solver.

AthenaK includes a built-in \emph{first-order flux correction} (FOFC) mechanism
\citep[see][]{Lemaster2009,Stone2020}.  When FOFC is active, a cell is flagged
as ``troubled'' if the PLM update would produce a non-physical state (negative
density or pressure).  The flux across every face of a flagged cell is then
recomputed using first-order (piecewise-constant) reconstruction and re-applied
before the state is finalised.  This is a local, post-hoc correction that
preserves conservation while preventing the solver from crashing.

FOFC requires at least three ghost cells so that the stencil for the
first-order recomputation is always available.  The relevant input parameters
are:

\begin{center}
\begin{tabular}{@{}lll@{}}
\toprule
Parameter & Block & Required value \\
\midrule
\texttt{nghost}  & \texttt{<mesh>}  & $\geq 3$ \\
\texttt{fofc}    & \texttt{<hydro>} / \texttt{<mhd>} & \texttt{true} \\
\bottomrule
\end{tabular}
\end{center}

Both are set in the canonical \texttt{sn\_drive\_box.athinput} input file.
Without FOFC, the simulation remains stable provided the injected radial
velocity $v_r$ is subsonic throughout the injection sphere.  With FOFC
enabled, transonic and mildly supersonic injection events are handled
gracefully at the cost of a small fraction of first-order cells per event.

\section{Conservation Properties}

\paragraph{Energy:} The total energy injected per event is
\begin{equation}
  \Delta E_\mathrm{tot} = E_\mathrm{inj} + \tfrac{1}{2}\,\bar{\rho}\,v_r^2\,V\,.
\end{equation}
The first term is the thermal contribution; the second is the kinetic energy
associated with the radial momentum.  Both are exactly conserved to
floating-point precision.

\paragraph{Linear momentum:} Because the radial velocity field is spherically
symmetric about the event centre, the net injected linear momentum vanishes by
symmetry:
\begin{equation}
  \sum_{c \,\in\, \mathcal{S}} \bar{\rho}\,v_r\,\hat{\boldsymbol{r}}_c\,\Delta V_c = \mathbf{0}\,.
\end{equation}
On a discrete Cartesian mesh this cancellation is exact for injection spheres
centred on a cell centre, and approximate (to machine precision) otherwise.

\paragraph{Mass:} No mass is injected.  The density field is unmodified by the
source term.

\section{Subgrid SN Remnant Model}
\label{sec:subgrid}

When \texttt{sn\_subgrid = true}, the per-event thermal energy and radial
momentum are no longer taken from the fixed input parameters
\texttt{sn\_energy}/\texttt{sn\_momentum}.  Instead, they are computed from
the fitting formulae of Martizzi, Faucher-Gigu\`ere \& Quataert (2015,
hereafter MFQ15) for the evolution of an SN remnant in a uniform or turbulent
ambient medium.  The independent-injection mode
(\texttt{sn\_subgrid = false}) remains the default and is unaffected.

\subsection{Physical Scales}

The fiducial SN energy and ejecta momentum are:
\begin{equation}
  E_\mathrm{th,0} = 0.717\,E_\mathrm{SN}\,,
  \qquad
  P_0 = \sqrt{2\,M_\mathrm{ej}\,E_\mathrm{SN}}\,,
\end{equation}
where $E_\mathrm{SN}$ and $M_\mathrm{ej}$ are user-specified (defaults
$10^{51}\,\mathrm{erg}$ and $3\,M_\odot$).  For these defaults
$P_0 \approx 3.45 \times 10^{42}\;\mathrm{g\,cm\,s^{-1}}$.

\subsection{Fitting Formulae}

Five characteristic radii (in pc) and one power-law index $\alpha$ are
evaluated from the local mean hydrogen number density $n_\mathrm{H}$ (in
$\mathrm{cm}^{-3}$, computed from $\bar{\rho}$ in the injection sphere) and a
fixed metallicity $Z/Z_\odot$ (floored at 0.01):
\begin{equation}
  X = X_\star\,\Bigl(\frac{n_\mathrm{H}}{100}\Bigr)^{a_X}
      \Bigl(\frac{Z}{Z_\odot}\Bigr)^{b_X}\,.
\end{equation}
Two fit sets are available, selected by \texttt{sn\_subgrid\_type}:

\paragraph{Type 0 --- Homogeneous medium (MFQ15 Eq.~11):}
\begin{center}
\begin{tabular}{@{}lccccc@{}}
\toprule
 & $\alpha$ & $R_c$ [pc] & $R_r$ [pc] & $R_0$ [pc] & $R_b$ [pc] \\
\midrule
$X_\star$ & $-7.8$ & 3.0 & 5.5 & 0.97 & 4.0 \\
$a_X$ ($n_\mathrm{H}$) & 0.050 & $-0.42$ & $-0.40$ & $-0.33$ & $-0.43$ \\
$b_X$ ($Z$) & 0.030 & $-0.082$ & $-0.074$ & 0.046 & $-0.077$ \\
\bottomrule
\end{tabular}
\end{center}

\paragraph{Type 1 --- Inhomogeneous medium, $\mathcal{M} = 30$ (MFQ15
Eq.~12):}
\begin{center}
\begin{tabular}{@{}lccccc@{}}
\toprule
 & $\alpha$ & $R_c$ [pc] & $R_r$ [pc] & $R_0$ [pc] & $R_b$ [pc] \\
\midrule
$X_\star$ & $-11$ & 6.3 & 9.2 & 2.4 & 8.0 \\
$a_X$ ($n_\mathrm{H}$) & 0.070 & $-0.42$ & $-0.44$ & $-0.35$ & $-0.46$ \\
$b_X$ ($Z$) & 0.114 & $-0.050$ & $-0.067$ & 0.021 & $-0.058$ \\
\bottomrule
\end{tabular}
\end{center}

\subsection{Per-Event Thermal Energy}

The retained thermal energy is evaluated as a piecewise function of the
effective injection radius $R = r_\mathrm{inj,eff}$ (converted to pc):
\begin{equation}\label{eq:Eth}
  E_\mathrm{th}(R) =
  \begin{cases}
    E_\mathrm{th,0}                                         & R \leq R_c
      \quad\text{(Sedov--Taylor)}\,, \\[4pt]
    E_\mathrm{th,0}\,(R / R_c)^{\alpha}                     & R_c < R \leq R_r
      \quad\text{(cooling)}\,, \\[4pt]
    E_\mathrm{th,0}\,(R_r / R_c)^{\alpha}                   & R > R_r
      \quad\text{(floor)}\,.
  \end{cases}
\end{equation}

\subsection{Per-Event Radial Momentum}

The radial momentum of the shell is:
\begin{equation}\label{eq:Prad}
  P_\mathrm{rad}(R) =
  \begin{cases}
    P_0\,(R / R_0)^{3/2}       & R \leq R_b
      \quad\text{(growing)}\,, \\[4pt]
    P_0\,(R_b / R_0)^{3/2}     & R > R_b
      \quad\text{(terminal)}\,.
  \end{cases}
\end{equation}

\subsection{Algorithm Integration}

The subgrid evaluation is inserted between the density-averaging pass
(Pass~2) and the injection pass (Pass~3) of Section~4.  The steps are:
\begin{enumerate}
  \item Convert $\bar{\rho}$ (code) $\to$ $n_\mathrm{H}$ (cm$^{-3}$) using
    the code density and mean molecular weight from the \texttt{<units>} block.
  \item Evaluate the five radii and $\alpha$ from the selected fit set.
  \item Compute $E_\mathrm{th}(R)$ and $P_\mathrm{rad}(R)$ via
    equations~\eqref{eq:Eth}--\eqref{eq:Prad}.
  \item Convert both to code units and use them in place of the fixed
    \texttt{sn\_energy}/\texttt{sn\_momentum} for this event.
\end{enumerate}
The density-averaging pass is always executed when \texttt{sn\_subgrid} is
active (it is needed to compute $n_\mathrm{H}$), even if the independent
\texttt{sn\_momentum} parameter is zero.

\subsection{Note on the RAMSES Implementation}

The RAMSES \texttt{blast\_wave\_feedback} routine (Beattie et al., in prep.)
uses the same MFQ15 fits, but the $n_\mathrm{H}$ and $Z$ exponents for
$\alpha$ appear to be \emph{swapped} relative to the published table
($a_\alpha = 0.024$, $b_\alpha = 0.050$ in RAMSES versus $a_\alpha = 0.050$,
$b_\alpha = 0.030$ in MFQ15 Eq.~11).  The AthenaK implementation uses the
paper values.

\section{Input Parameters}

All SN driving parameters are read from the \texttt{<hydro>} (or
\texttt{<mhd>}) input block.  Table~\ref{tab:params} lists the full set.

\begin{table}[h]
\centering
\caption{Supernova driving input parameters.}\label{tab:params}
\begin{tabular}{@{}llll@{}}
\toprule
Parameter & Type & Default & Description \\
\midrule
\multicolumn{4}{@{}l}{\emph{Core driving parameters}} \\
\texttt{sn\_driving}       & bool   & \texttt{false} & Enable SN driving \\
\texttt{sn\_rate}          & Real   & (required)      & Events per unit code time \\
\texttt{sn\_energy}        & Real   & ---             & Thermal energy per event (code units) \\
\texttt{sn\_energy\_cgs}   & Real   & $10^{51}$       & Thermal energy per event (erg) \\
\texttt{sn\_momentum}      & Real   & ---             & Radial momentum per event (code units) \\
\texttt{sn\_momentum\_cgs} & Real   & ---             & Radial momentum per event (g\,cm\,s$^{-1}$) \\
\texttt{sn\_rinj}          & Real   & ---             & Injection radius (code units) \\
\texttt{sn\_radius\_pc}    & Real   & 8.0             & Injection radius (pc) \\
\texttt{sn\_zmin}          & Real   & $x_{3,\min}$    & Lower $z$-band edge (code units) \\
\texttt{sn\_zmax}          & Real   & $x_{3,\max}$    & Upper $z$-band edge (code units) \\
\texttt{sn\_zmin\_pc}      & Real   & ---             & Lower $z$-band edge (pc) \\
\texttt{sn\_zmax\_pc}      & Real   & ---             & Upper $z$-band edge (pc) \\
\texttt{sn\_seed}          & int    & 12345           & RNG seed \\
\texttt{sn\_log\_events}   & bool   & \texttt{true}   & Write event log file \\
\texttt{sn\_log\_file}     & string & \texttt{<basename>.sn\_events.dat} & Event log filename \\
\midrule
\multicolumn{4}{@{}l}{\emph{Subgrid model parameters (Section~\ref{sec:subgrid})}} \\
\texttt{sn\_subgrid}       & bool   & \texttt{false}  & Enable subgrid $E_\mathrm{th}$/$P_\mathrm{rad}$ model \\
\texttt{sn\_subgrid\_type} & int    & 0               & Fit set: 0 = homogeneous, 1 = inhomogeneous \\
\texttt{sn\_esn\_cgs}      & Real   & $10^{51}$       & Total SN energy [erg] \\
\texttt{sn\_mej\_msun}     & Real   & 3.0             & Ejecta mass [$M_\odot$] \\
\texttt{sn\_metallicity}   & Real   & 1.0             & Metallicity [$Z/Z_\odot$] \\
\bottomrule
\end{tabular}
\end{table}

For \texttt{sn\_energy} and \texttt{sn\_momentum}, code-unit and CGS variants
are mutually exclusive; the code-unit version takes precedence if both are
specified.  If neither momentum parameter is provided, $P_\mathrm{inj} = 0$ and
the module reverts to thermal-only injection (the original behaviour).

When \texttt{sn\_subgrid = true}, the values of \texttt{sn\_energy} and
\texttt{sn\_momentum} are ignored; $E_\mathrm{th}$ and $P_\mathrm{rad}$ are
computed per-event from the fitting formulae.  The event log file records the
computed injection values for each event.

The CGS momentum conversion uses:
\begin{equation}
  P_\mathrm{code} = \frac{P_\mathrm{cgs}}{M_\mathrm{cgs}\,V_\mathrm{cgs}}\,,
\end{equation}
where $M_\mathrm{cgs}$ and $V_\mathrm{cgs}$ are the code mass and velocity
scales in CGS.

\section{Relation to the RAMSES Implementation}

The momentum injection follows the same physical approach as the
\texttt{blast\_wave\_feedback} routine in RAMSES (Beattie et al., in prep.),
with one simplification:

\begin{enumerate}
  \item \textbf{No ejecta mass.}  The RAMSES implementation overwrites the
    density within the injection sphere with $\rho_\mathrm{ej} + \bar{\rho}$,
    where $\rho_\mathrm{ej} = M_\mathrm{ej} / V_\mathrm{inj}$ depends on the
    event type (SNII, SNIa, NSM).  The AthenaK implementation does not modify
    the density field; momentum is injected into the existing gas.
\end{enumerate}

The radial velocity in RAMSES is:
\begin{equation}
  v_r^\mathrm{RAMSES} = \frac{3\,P_\mathrm{rad}}{4\pi\,(\rho_\mathrm{ej} + \bar{\rho})\,r_\mathrm{max}^3}\,,
\end{equation}
which reduces to equation~\eqref{eq:vri} in the limit $\rho_\mathrm{ej} \to 0$
and $V \to \frac{4}{3}\pi\,r_\mathrm{inj}^3$ (continuous sphere).

\section{Implementation Files}

\begin{description}
  \item[\texttt{src/srcterms/srcterms.hpp}] Class declaration; \texttt{sn\_pinj}
    data member.
  \item[\texttt{src/srcterms/srcterms.cpp}] Constructor (parameter parsing) and
    \texttt{SupernovaDriving()} (three-pass injection algorithm).
  \item[\texttt{inputs/hydro/sn\_drive\_box.athinput}] Example input file for a
    periodic box with stochastic SN driving.
\end{description}

\end{document}
